\documentclass[11pt, oneside]{article}   	% use "amsart" instead of "article" for AMSLaTeX format
\usepackage{geometry}                		% See geometry.pdf to learn the layout options. There are lots.
\geometry{letterpaper}                   		% ... or a4paper or a5paper or ... 
%\geometry{landscape}                		% Activate for rotated page geometry
\usepackage[parfill]{parskip}    		% Activate to begin paragraphs with an empty line rather than an indent
\usepackage{graphicx}				% Use pdf, png, jpg, or eps§ with pdflatex; use eps in DVI mode
								% TeX will automatically convert eps --> pdf in pdflatex		
\usepackage{amssymb}

%SetFonts

%SetFonts


\title{Math F113X: Homework Set 1}
%\author{The Author}
\date{}							% Activate to display a given date or no date

\begin{document}
\maketitle
%\section{}
%\subsection{}

%Homework assignment 1 is:

\fbox{\parbox{\textwidth}{
\begin{itemize}
\item Answer the following problems from the Voting Theory section (p.53):
\begin{quote}\# 1, 3, 6, 8, 10, 17, 20, Problem A\end{quote}
\end{itemize}
}} \\

{\bf Problem A:}

Consider the following preference schedule for an election with 4 candidates and 100 voters.

\begin{center}
\begin{tabular}{c||c|c|c|c|c|}
&24&10&21&27&18\\
\hline\hline
1st&A&B&D&C&B\\
2nd&C&A&A&D&A\\
3rd&D&C&C&A&C\\
4th&B&D&B&B&D
\end{tabular}
\end{center}

If an IRV (ranked choice) election was held, the first round
would not produce a winner, but D would be eliminated. Knowing this,
determine the outcome of the {\bf second round}, showing your
work. Give the winner (if there is one), or say who is eliminated
next. \\

\vspace{1cm}

Remember to write up your homework solutions according to the homework writeup guidelines. \\

After completing each problem, you should check your work and your answers against the solutions provided in Canvas!\\

Confirm your work is correct or fix your errors before you submit it. \\

Homework is graded using the following rubric for each problem (or problem part):

\begin{description}
\item[2 points:] You provided a complete answer, with supporting work, written up clearly
\item[1 point:] Some attempt at a solution, but incomplete writeup / unclear / illegible / no answer
\item[1 point:] Only an answer, with no supporting work 
\item[0 points:] Missing.
\end{description}



\end{document}  
%%% Local Variables:
%%% mode: LaTeX
%%% TeX-master: t
%%% End:
